\section{Berechnung der Verlustleistungen}
\label{app:verlustrechnung}

Dieser Anhang dokumentiert die mathematischen Zwischenschritte der in Abschnitt \ref{sec:verlustanalyse} durchgeführten Verlustleistungsberechnungen. Alle Rechnungen basieren auf den identifizierten Worst-Case-Szenarien ($I_{\mathrm{Block}} = 42{,}4\,\mathrm{A}$, $v_{\mathrm{i}} = 42\,\mathrm{V}$).

\subsection{Leitverluste MOSFET (High-Side und Low-Side)}
Es wird die Synchrongleichrichtung bei einer Sperrschichttemperatur von $T_{\mathrm{J,design}} = 100^\circ\mathrm{C}$ ($R_{\mathrm{DS(on)}} = 0{,}0048\,\Omega$) abgebildet. 

\textbf{Worst-Case High-Side-Schalter} (Maximale Aussteuerung $d=1$):
\begin{align}
    I_{\mathrm{HS,RMS}}^2 &= (42{,}4\,\mathrm{A})^2 \cdot \frac{1}{3} \cdot 1 = 599{,}25\,\mathrm{A}^2 \\
    P_{\mathrm{HS,c}} &= 599{,}25\,\mathrm{A}^2 \cdot 0{,}0048\,\Omega \approx 2{,}88\,\mathrm{W}
\end{align}

\textbf{Worst-Case Low-Side-Schalter} (Maximaler Freilauf bei Anlauf $d=0$):
\begin{align}
    I_{\mathrm{LS,RMS}}^2 &= (42{,}4\,\mathrm{A})^2 \cdot \frac{1}{3} \cdot (2 - 0) = 1797{,}76\,\mathrm{A}^2 \cdot \frac{1}{3} \cdot 2 \approx 1138{,}58\,\mathrm{A}^2 \\
    P_{\mathrm{LS,c}} &= 1138{,}58\,\mathrm{A}^2 \cdot 0{,}0048\,\Omega \approx 5{,}75\,\mathrm{W}
\end{align}

\subsection{Schaltverluste (High-Side)}
Bestimmung der Gateströme in der Miller-Phase ($v_{\mathrm{dr}} = 12\,\mathrm{V}$, $v_{\mathrm{GS,Mi}} = 5{,}0\,\mathrm{V}$, $R_{\mathrm{G}} = 10\,\Omega$):
\begin{align}
    i_{\mathrm{G,on}} &= \frac{12\,\mathrm{V} - 5{,}0\,\mathrm{V}}{10\,\Omega} = 0{,}7\,\mathrm{A} \\
    i_{\mathrm{G,off}} &= \frac{5{,}0\,\mathrm{V} - 0\,\mathrm{V}}{10\,\Omega} = 0{,}5\,\mathrm{A}
\end{align}

Berechnung der Umschaltdauern mit der Umschaltladung $Q_{\mathrm{sw}} = 26\,\mathrm{nC}$:
\begin{align}
    t_{\mathrm{on}} &= \frac{26\,\mathrm{nC}}{0{,}7\,\mathrm{A}} \approx 37{,}14\,\mathrm{ns} \\
    t_{\mathrm{off}} &= \frac{26\,\mathrm{nC}}{0{,}5\,\mathrm{A}} = 52{,}0\,\mathrm{ns}
\end{align}

Berechnung der Schaltenergien ($v_{\mathrm{i}} = 42\,\mathrm{V}$, $I_{\mathrm{Block}} = 42{,}4\,\mathrm{A}$):
\begin{align}
    E_{\mathrm{MOS,on}} &= \frac{1}{2} \cdot 42{,}4\,\mathrm{A} \cdot 42\,\mathrm{V} \cdot 37{,}14\,\mathrm{ns} \approx 33{,}07\,\mu\mathrm{J} \\
    E_{\mathrm{on,rr}} &= 332\,\mathrm{nC} \cdot 42\,\mathrm{V} \approx 13{,}94\,\mu\mathrm{J} \quad \text{(Reverse Recovery)} \\
    E_{\mathrm{on,total}} &= 33{,}07\,\mu\mathrm{J} + 13{,}94\,\mu\mathrm{J} = 47{,}01\,\mu\mathrm{J} \\
    E_{\mathrm{MOS,off}} &= \frac{1}{2} \cdot 42{,}4\,\mathrm{A} \cdot 42\,\mathrm{V} \cdot 52{,}0\,\mathrm{ns} \approx 46{,}30\,\mu\mathrm{J}
\end{align}

Dynamische Verlustleistung ($f_{\mathrm{sw}} = 20\,\mathrm{kHz}$ bei aktiven $120^\circ$ pro Periode):
\begin{equation}
    P_{\mathrm{MOS,sw}} = (47{,}01\,\mu\mathrm{J} + 46{,}30\,\mu\mathrm{J}) \cdot 20000\,\mathrm{Hz} \cdot \frac{1}{3} \approx 0{,}62\,\mathrm{W}
\end{equation}

\subsection{Leitverluste Freilaufdiode (Totzeitverluste)}
Die Diode leitet aus Sicherheitsgründen nur während der Treiber-Totzeit $t_{\mathrm{dt}} = 520\,\mathrm{ns}$.
\begin{align}
    d_{\mathrm{D}} &= 2 \cdot 520\,\mathrm{ns} \cdot 20\,\mathrm{kHz} = 0{,}0208 \\
    \overline{i}_{\mathrm{D}} &= 42{,}4\,\mathrm{A} \cdot 0{,}0208 \cdot \frac{1}{3} \approx 0{,}294\,\mathrm{A} \\
    i_{\mathrm{D,RMS}}^2 &= (42{,}4\,\mathrm{A})^2 \cdot 0{,}0208 \cdot \frac{1}{3} \approx 12{,}46\,\mathrm{A}^2
\end{align}

Leitverlustleistung mit den grafisch ermittelten Parametern ($v_{\mathrm{D,0}} = 0{,}58\,\mathrm{V}$, $r_{\mathrm{D}} = 0{,}008\,\Omega$):
\begin{equation}
    P_{\mathrm{D,c}} = 0{,}58\,\mathrm{V} \cdot 0{,}294\,\mathrm{A} + 0{,}008\,\Omega \cdot 12{,}46\,\mathrm{A}^2 = 0{,}171\,\mathrm{W} + 0{,}100\,\mathrm{W} \approx 0{,}27\,\mathrm{W}
\end{equation}